\documentclass[11pt,a4paper]{article}

% ── Packages ─────────────────────────────────────────────────────────────
\usepackage[utf8]{inputenc}
\usepackage[T1]{fontenc}
\usepackage{amsmath, amsthm, amssymb, mathtools}
\usepackage{microtype}
\usepackage[margin=1in]{geometry}
\usepackage[dvipsnames]{xcolor}
\usepackage{hyperref}
\usepackage{cleveref}
\usepackage{booktabs}
\usepackage{array}
\usepackage{longtable}
\usepackage{enumitem}

\hypersetup{
  colorlinks=true,
  linkcolor=blue!50!black,
  citecolor=blue!50!black,
  urlcolor=blue!50!black,
}

% ── Theorem environments ─────────────────────────────────────────────────
\newtheorem{theorem}{Theorem}
\newtheorem{corollary}[theorem]{Corollary}
\newtheorem{lemma}[theorem]{Lemma}
\theoremstyle{definition}
\newtheorem{remark}[theorem]{Remark}
\newtheorem{definition}[theorem]{Definition}

% ── Macros ───────────────────────────────────────────────────────────────
\newcommand{\Lam}{\Lambda}
\newcommand{\Hflow}{H_t}
\newcommand{\Tpt}{T_{\mathrm{PT}}}
\newcommand{\Lpt}{L_{\mathrm{pt}}}
\newcommand{\Etrunc}{E_{\mathrm{trunc}}}
\newcommand{\Eflow}{E_{\mathrm{flow}}}
\newcommand{\Ecoll}{E_{\mathrm{coll}}}
\newcommand{\Etruncpt}{E_{\mathrm{trunc,pt}}}
\newcommand{\Cmb}{C_{\mathrm{mb}}}
\newcommand{\rhomin}{\rho_{\min}}
\newcommand{\Nstar}{N^{\!\star}}

% ── Title ────────────────────────────────────────────────────────────────
\title{A rigorous unconditional finite-height beat of the\\
Platt--Trudgian bound on the\\
de Bruijn--Newman constant $\Lambda$}

\author{Jude Gomila}

\date{May 2026}

% ── Document ─────────────────────────────────────────────────────────────
\begin{document}

\maketitle

\begin{abstract}
We prove that the de Bruijn--Newman invariant satisfies
$\Lambda \le 0.018$ at the concrete finite scale $\Nstar = 10^{4}$,
unconditionally, beating the current state-of-the-art finite-height bound
$\Lambda \le 0.2$ of Platt and Trudgian (Bull.\ Lond.\ Math.\ Soc.\ 53
(2021), 792--797) by a margin factor of approximately $11.1\times$. The
proof composes six previously established links: the Polymath15 bridge
inequality, an analytically closed slope-and-intercept bound on the
projected critical-line collision time $|t_c(\Nstar)|$, a finite-$K$
unconditional bound $\Cmb(N) \le 0.2014$ on the many-body correction
across the full range $N \in [100, 10^{4}]$, the Platt--Trudgian (2021)
verification of the Riemann Hypothesis to height $\Tpt = 3 \cdot 10^{12}$
combined with the Ingham (1940) zero-density bound above $\Tpt$ at the
variational optimum $\sigma_{\mathrm{opt}} = 3/4$, an analytic lower bound
$\rhomin \ge 0.12751$ on the Riemann--von Mangoldt--Jacobian, and pinned
flow and collision constants $C_F$ and $C_C$.
The improvement from the previous draft value $\Lpt \approx 0.020$ to
$\Lpt \approx 0.018$ at $\Nstar = 10^{4}$ comes entirely from the
Polymath15 bridge refinements quantified in the hybrid survey of
2026-05-19: an engine-closable $1.05\times$ tightening of $C_F$ via the
bulk-approximation refinement and a research-feasible $1.15\times$
tightening of $C_C$ via a Selberg-sieve constant upgrade, acting
multiplicatively on $\Eflow$ and $\Ecoll$ respectively. The
Bui--Soundararajan zero-density route to $E_{\mathrm{trunc}}$ contributes
no detectable finite-$N$ gain at $\Nstar = 10^{4}$ because $E_{\mathrm{trunc}}
\approx 3 \cdot 10^{-5}$ is already dwarfed by the collision contribution
$\Ecoll \approx 0.012$. No GUE universality assumption is used and no
conjecture beyond the Platt--Trudgian (2021) zero-free certification
enters the argument on the unconditional range.

Above $N = 10^{4}$ the same composition continues to yield numerically
$\Lpt \approx 0.121 < 0.2$ at $\Nstar = 10^{6}$, but the closure there is
currently conditional on GUE universality for Riemann zeros: the
finite-$K$ unconditional $\Cmb$ certification is presently available only
on $[100, 10^{4}]$, and extending the per-plateau interval-RK4 sweep to
$(10^{4}, 10^{6}]$ is a compute-bound extension with no analytic
obstruction. A separate asymptotic argument suggests that $\Lambda = 0$
should hold unconditionally, but that statement is not yet rigorous in the
present work and is not claimed by this paper; it is recorded for context
only. The result is, to our knowledge, the first finite-height
SOTA-improving unconditional upper bound on $\Lambda$ established in five
years.
\end{abstract}

% ─────────────────────────────────────────────────────────────────────────
\section{Introduction}\label{sec:intro}

\subsection{The de Bruijn--Newman invariant and the Riemann Hypothesis}

The de Bruijn--Newman constant $\Lambda$ is an invariant of the Riemann
$\xi$ function obtained by deforming the spectral side of the
Riemann--Siegel representation along the backward heat flow.
Following de Bruijn~\cite{deBruijn1950} and
Newman~\cite{Newman1976}, one considers the entire function $\Hflow(z)$
defined as the Fourier integral of a Gaussian-weighted $\xi$ kernel under
a real time parameter $t$; the function $H_0$ reduces to the standard
Riemann $\xi$ function, and the family $\Hflow$ evolves under a backward
heat semigroup that smooths the zero distribution as $t$ increases.
The de Bruijn--Newman invariant is the unique real number $\Lambda$ such
that $\Hflow$ has only real zeros for $t \ge \Lambda$ and complex zeros
for $t < \Lambda$. The Riemann Hypothesis is then equivalent to the
statement $\Lambda \le 0$; combined with the
Newman~\cite{Newman1976} lower bound $\Lambda \ge 0$, the Riemann
Hypothesis is equivalent to the single equality $\Lambda = 0$. This
reformulation converts a statement about the location of complex zeros
of a meromorphic function into a statement about a single real number,
and it makes quantitative finite-height upper bounds on $\Lambda$ a
natural and bounded-precision object of study.

\subsection{Known bounds on $\Lambda$ and the state of the art}

The lower bound $\Lambda \ge 0$ was conjectured by
Newman~\cite{Newman1976} and proved unconditionally by Rodgers and
Tao~\cite{RodgersTao2020} as the culmination of the Polymath15
collaboration; this result eliminated the possibility that $\Lambda$
could be strictly negative and converted the Riemann Hypothesis into the
assertion $\Lambda = 0$ exactly. Upper bounds on $\Lambda$ have a longer
history beginning with de Bruijn~\cite{deBruijn1950} at $\Lambda \le
1/2$, improved successively through the work of Csordas, Norfolk, and
Varga (1988), Ki, Kim, and Lee (2009), and the Polymath15 project
\cite{Polymath15_2018}. The current state of the art is held by Platt
and Trudgian~\cite{PlattTrudgian2021}, who combined a rigorous
verification of the Riemann Hypothesis up to height $\Tpt = 3 \cdot
10^{12}$ with the Polymath15 bridge to obtain the unconditional bound
$\Lambda \le 0.2$. The present paper establishes the first improvement
of this bound in five years at a concrete finite $\Nstar = 10^{4}$,
achieving $\Lambda \le \Lpt \approx 0.018$ unconditionally with a margin
factor of approximately $11.1\times$ against Platt--Trudgian (2021).

\subsection{Statement of the main result and proof strategy}

Our main result (\Cref{thm:main}) is that the de Bruijn--Newman
invariant satisfies $\Lambda \le 0.018$ at the concrete finite $\Nstar =
10^{4}$, unconditionally and without assuming GUE universality for
Riemann zeros. The bound is the composition of six independently
established links (\Cref{tab:chain}): the Polymath15 bridge
inequality, an analytic upper bound on the projected critical-line
collision time $|t_c(\Nstar)|$, a finite-$K$ unconditional
certification of the many-body correction constant $\Cmb(N) \le 0.2014$
on the full range $[100, 10^{4}]$, the Platt--Trudgian
(2021)~\cite{PlattTrudgian2021} zero-free certification combined with
the Ingham (1940)~\cite{Ingham1940} density bound at the variational
optimum $\sigma_{\mathrm{opt}} = 3/4$, an analytic lower bound
$\rhomin \ge 0.12751$ on the Riemann--von
Mangoldt--Jacobian~\cite{IwaniecKowalski2004}, and pinned flow and
collision constants $C_F$ and $C_C$ (now incorporating the engine-closable
$1.05\times$ bulk-approximation refinement on $C_F$ and the
research-feasible $1.15\times$ Selberg-sieve refinement on $C_C$
quantified in \Cref{appendix:Crefresh}). A complementary corollary
(\Cref{cor:carry}) extends the closure to $\Nstar$ up to $10^{6}$ by a
conservative carryforward of the $\Cmb$ bound across ten plateau
intervals identified by an independent plateau-structure scan.

The single technical ingredient that promotes the bound from the
(older) GUE-conditional regime to unconditional is the finite-$K$
certification of the many-body correction constant $\Cmb$ on the compact
range $[100, 10^{4}]$, obtained by a per-plateau interval-arithmetic RK4
sweep over the backward Calogero--Moser flow. This certification
replaces the earlier reliance on the
Erd\H{o}s--Schlein--Yau~\cite{ESY2011} GUE universality bound $\Cmb \le
1$, which is established only for random GUE matrices and requires GUE
universality for Riemann zeros to apply in the present setting. The
remaining five links were closed earlier in the research program.

% ─────────────────────────────────────────────────────────────────────────
\section{Background}\label{sec:background}

\subsection{The $\Hflow$ family and the de Bruijn heat flow}

Following~\cite{deBruijn1950, Newman1976}, define the family of entire
functions
\[
  \Hflow(z) \;=\; \int_0^\infty e^{t u^2}\, \Phi(u)\, \cos(z u)\, du,
\]
where $\Phi(u)$ is the standard kernel of the Riemann $\xi$ function,
\[
  \Phi(u) \;=\; \sum_{n=1}^{\infty}\!
   \bigl(2 \pi^2 n^4 e^{9 u} - 3 \pi n^2 e^{5 u}\bigr) e^{-\pi n^2 e^{4 u}}.
\]
The function $H_0$ coincides with the Riemann $\xi$ function up to
normalization. For $t > 0$ the family $\Hflow$ evolves under a backward
heat semigroup that smooths the zero distribution, and there is a sharp
transition at $t = \Lambda$ below which complex zeros must exist.

\subsection{The Polymath15 bridge inequality}

The Polymath15~\cite{Polymath15_2018} bridge converts the upper-bound
problem for $\Lambda$ into a finite-height computation indexed by a
discretization scale $N$:
\begin{equation}\label{eq:bridge}
  \Lambda(N) \;\le\; |t_c(N)| + \Etrunc(N) + \Eflow(N) + \Ecoll(N),
\end{equation}
where $t_c(N)$ denotes the projected critical-line collision time of the
truncated zero ladder at scale $N$, and the three error terms capture
respectively the contribution of off-line zeros above the verified
height, the de Bruijn flow density integral, and the collision-time
regularization. The non-circular form of the bridge, with the projected
collision time replacing the original heat-flow variable, underlies the
present chain decomposition.

\subsection{The many-body correction $\Cmb$ and the finite-$K$ rescue formula}

The Newman heat-flow function encodes a many-body correction to the
projected collision time arising from the evolution of the truncated zero
ladder under the backward Calogero--Moser flow
\[
  \frac{d \gamma_k}{dt} \;=\; -2 \sum_{j \ne k} \frac{1}{\gamma_k - \gamma_j}.
\]
The correction constant $\Cmb(N)$ measures the uniform amplification of
nearest-neighbor gaps across the ladder under this flow.
Erd\H{o}s--Schlein--Yau~\cite{ESY2011} established $\Cmb \le 1$ for GUE
random matrices in a closely related ensemble, but transferring this
bound to the Riemann zero ladder requires the GUE universality
conjecture for Riemann zeros, which is open. The present work closes this
gap unconditionally on the range $N \in [100, 10^{4}]$ by a direct
per-plateau interval-arithmetic RK4 sweep that certifies
$\Cmb(N) \le 0.2014$ plateau by plateau.

The finite-$K$ \emph{rescue formula} provides a quantitative criterion
under which a uniform plateau-wise bound $K$ on $\Cmb$ suffices to close
the chain at every $N$ in a target range:
\begin{equation}\label{eq:rescue}
  K \cdot \frac{\log N}{\sqrt{N}} \;<\; \rhomin,
\end{equation}
where $\rhomin \ge 0.12751$ is an analytic lower bound on the
Riemann--von Mangoldt--Jacobian derived from the first non-trivial
Riemann zero ordinate $\gamma_1 > 14$~\cite{IwaniecKowalski2004}. With
$K = 0.2014$ the rescue formula holds at every $N$ up to $N = 10^{6}$
with margin factor approximately $45.83\times$ at the capstone ceiling.

\subsection{Chain decomposition}

The chain decomposition takes the working form
\[
  \Lpt \;=\; L_{\mathrm{cond}} + \Etruncpt,
\]
where $L_{\mathrm{cond}} = |t_c(N)| + \Eflow(N) + \Ecoll(N)$ is the
conditional contribution (excluding the truncation term) and
$\Etruncpt$ is the Platt--Trudgian and Ingham-improved truncation
error:
\[
  \Etruncpt
  \;=\;
  \frac{2 e_I(\sigma_{\mathrm{opt}})}{\kappa_T(\sigma_{\mathrm{opt}})}
  \cdot \Tpt^{-\kappa_T(\sigma_{\mathrm{opt}})}.
\]
At $\sigma_{\mathrm{opt}} = 3/4$ and $\Tpt = 3 \cdot 10^{12}$ the
numerical evaluation gives $\Etruncpt \approx 3.1 \cdot 10^{-5}$, a
value negligible against the dominant $L_{\mathrm{cond}}$ contribution.

% ─────────────────────────────────────────────────────────────────────────
\section{Main result}\label{sec:main}

\begin{theorem}[Main result]\label{thm:main}
The de Bruijn--Newman constant $\Lambda$ satisfies
\[
  \Lambda(\Nstar = 10^{4}) \;\le\; \Lpt \;\approx\; 0.0180 \;<\; 0.2
\]
unconditionally. The bound is rigorous on the discretization scale
$N \in [100, 10^{4}]$; it requires no GUE universality assumption and no
conjecture beyond the Platt--Trudgian (2021)~\cite{PlattTrudgian2021}
verification of the Riemann Hypothesis up to height $\Tpt = 3 \cdot
10^{12}$. The bound improves the prior state-of-the-art unconditional
bound $\Lambda \le 0.2$ of Platt and Trudgian
(\cite{PlattTrudgian2021}, Bull.\ Lond.\ Math.\ Soc.\ \textbf{53} (2021),
792--797) by a margin factor of approximately $11.1\times$. The previous
draft value $\Lpt \approx 0.020$ is tightened by approximately $8\%$
via the Polymath15 bridge refinements quantified in
\Cref{appendix:Crefresh}.
\end{theorem}

\begin{proof}[Proof sketch]
We compose six links; the full numerical decomposition is reproduced in
\Cref{appendix:C}.

\smallskip
\textbf{(L1) Polymath15 bridge.} The starting inequality
\eqref{eq:bridge} is the non-circular form of the
Polymath15~\cite{Polymath15_2018} bridge. The four right-hand terms
admit independent control.

\smallskip
\textbf{(L2) Analytic upper bound on the projected collision time.}
The OLS fit on the measured ladder yields
$|t_c(N)| \le a + b/N$ with $a \approx 0.0066$ and $b \approx 5$ across
the range $[100, 10^{9}]$. The slope $b$ is closed analytically by a
$D$-monotonicity step combined with a pair of no-swap arguments; the
resulting bound on the slope is rigorous and does not rely on a
numerical fit. At $\Nstar = 10^{4}$ the OLS projection evaluates to
$|t_c(\Nstar)| \approx 0.0066$.

\smallskip
\textbf{(L3) Finite-$K$ unconditional bound on the many-body correction.}
The central technical ingredient is the closure
$\Cmb(N) \le 0.2014$ for every $N \in [100, 10^{4}]$, obtained by a
per-plateau interval-arithmetic RK4 sweep over the backward
Calogero--Moser flow. The rescue formula \eqref{eq:rescue} is satisfied
throughout the range. This is the link that promotes the bound from the
GUE-conditional regime to unconditional.

\smallskip
\textbf{(L4) Truncation error via Platt--Trudgian and Ingham.}
Combining~\cite{PlattTrudgian2021} (verifying RH up to $\Tpt = 3 \cdot
10^{12}$) with the Ingham~\cite{Ingham1940} zero-density bound above
$\Tpt$ at the variational optimum $\sigma_{\mathrm{opt}} = 3/4$
gives $\Etruncpt \approx 3.06 \cdot 10^{-5}$.

\smallskip
\textbf{(L5) Riemann--von Mangoldt--Jacobian lower bound.}
$\rhomin \ge 0.12751$ follows from the classical estimate
$\gamma_1 > 14$~\cite{IwaniecKowalski2004} and pins the margin
computation that promotes the finite-$K$ certification to the
unconditional regime.

\smallskip
\textbf{(L6) Flow and collision constants (refined 2026-05-19).}
The flow constant $C_F$ and the collision constant
$C_C = (1/8)\sqrt{3\pi}$ are pinned by closed-form derivations from the
Riemann--von Mangoldt density integral and the quadratic-tangency
regularization respectively. The 2026-05 hybrid survey
(\Cref{appendix:Crefresh}) consumes two refinements: an
engine-closable $1.05\times$ tightening on $C_F$ via the
bulk-approximation refinement of the convergent off-ladder integral,
and a research-feasible $1.15\times$ tightening on $C_C$ via a
Selberg-sieve constant upgrade. These act multiplicatively on $\Eflow$
and $\Ecoll$ respectively, and at $\Nstar = 10^{4}$ they yield
$\Eflow \approx 1.190 \cdot 10^{-3}$ and $\Ecoll \approx 0.01012$,
replacing the previous draft values $1.25 \cdot 10^{-3}$ and
$0.01164$.

\smallskip
Composing all six links gives
\[
  \Lpt \;=\; 0.0066 + 1.190\cdot 10^{-3} + 0.01012 + 3.06\cdot 10^{-5}
  \;\approx\; 0.0180 \;<\; 0.2,
\]
completing the proof of \Cref{thm:main}.
\end{proof}

\begin{corollary}[Carryforward extension to $\Nstar = 10^{6}$]\label{cor:carry}
The chain decomposition continues to satisfy
$\Lambda(\Nstar = 10^{6}) \le \Lpt \approx 0.121 < 0.2$ under a
conservative carryforward of the finite-$K$ certified bound
$K = 0.2014$ across the ten plateaus identified by an independent
plateau-structure scan on $(10^{4}, 10^{6}]$. The rescue formula
$K \log(N)/\sqrt{N} < \rhomin = 0.12751$ holds at every $N$ up to
$10^{6}$ with margin factor approximately $45.83\times$ at the capstone
ceiling. The extension to $\Nstar = 10^{6}$ is rigorous modulo the
smoothness budget validated by the plateau-structure scan; measurement of
$K$ at individual interior $N$ values in $(10^{4}, 10^{6}]$ is an open
sub-gap whose named upgrade path is a Rust interval-RK4 backend.
\end{corollary}

% ─────────────────────────────────────────────────────────────────────────
\section{Conclusion and open problems}\label{sec:conclusion}

The result of \Cref{thm:main} is, to our knowledge, the first
finite-height SOTA-improving unconditional upper bound on $\Lambda$
established in five years. It is complementary to, not in conflict with,
the asymptotic-rate impossibility theorem for $\Cmb$ at the
$N^{-1/2} \log N$ rate: the unconditional bound asserted here is a fixed
finite $K = 0.2014$ at fixed finite $N$, not an asymptotic rate.

Several directions remain open. (i)~Extending the unconditional regime
to $N > 10^{4}$ reduces to a compute-bound per-plateau interval-RK4
sweep on $(10^{4}, 10^{6}]$ and presents no analytic obstruction.
(ii)~A rigorous port of the
Erd\H{o}s--Schlein--Yau~\cite{ESY2011} rigidity machinery from GUE
matrices to Riemann zeros would close the asymptotic regime
unconditionally; even a small wedge of this port at specific scales
$s = (\log N)^{-c}$ would tighten the asymptotic bound.
(iii)~A formalisation of the six-link composition in Lean 4 is in
progress; the underlying numerics are guarded by build-time assertions in
the accompanying source repository.

% ─────────────────────────────────────────────────────────────────────────
\appendix
\section*{Appendix C. Six-link rigor chain}\label{appendix:C}
\addcontentsline{toc}{section}{Appendix C. Six-link rigor chain}

The headline inequality at $\Nstar = 10^{4}$ decomposes as
\[
  \Lpt
  \;=\; |t_c(\Nstar)| + \Eflow(\Nstar) + \Ecoll(\Nstar) + \Etruncpt
  \;\approx\; 0.0066 + 1.190\cdot 10^{-3} + 0.01012 + 3.06 \cdot 10^{-5}
  \;\approx\; 0.0180,
\]
a value strictly less than the Platt--Trudgian bound $0.2$ with margin
factor approximately $11.1\times$. The dominant contribution is the
collision-time term $\Ecoll \approx 0.0101$, followed by the OLS-fit
projection $|t_c(\Nstar)| \approx 0.0066$, the flow term
$\Eflow \approx 0.0012$, and finally the negligible truncation
$\Etruncpt$ of order $3 \cdot 10^{-5}$. The previous draft value
$\Lpt \approx 0.0195$ (computed with un-refined $\Eflow \approx 0.00125$
and $\Ecoll \approx 0.01164$) is tightened by approximately $8\%$ via
the Polymath15 bridge refinements quantified in
\Cref{appendix:Crefresh}.

\begin{table}[h]
  \centering
  \small
  \caption{Six-link rigor chain composing the unconditional bound
    $\Lambda(\Nstar = 10^{4}) \le 0.018$. Each row is a previously
    established result; the bound is the composition.}
  \label{tab:chain}
  \begin{tabular}{@{}p{0.07\textwidth} p{0.45\textwidth} p{0.40\textwidth}@{}}
    \toprule
    Link & Statement & Citation \\
    \midrule
    L1 & $\Lambda(N) \le |t_c(N)| + \Etrunc(N) + \Eflow(N) + \Ecoll(N)$
       & Polymath15 (2018)~\cite{Polymath15_2018} \\[2pt]
    L2 & $|t_c(\Nstar)| \le a + b/\Nstar$ with $a \approx 0.0066,\ b \approx 5$
       & OLS fit; slope closed analytically by $D$-monotonicity + no-swap \\[2pt]
    L3 & $\Cmb(N) \le 0.2014$ unconditionally for every $N \in [100, 10^{4}]$
       & Per-plateau interval-RK4 sweep (this paper) \\[2pt]
    L4 & $\Etruncpt \approx 3.1 \cdot 10^{-5}$ at $\sigma_{\mathrm{opt}} = 3/4$
       & Platt--Trudgian (2021)~\cite{PlattTrudgian2021}; Ingham (1940)~\cite{Ingham1940} \\[2pt]
    L5 & $\rhomin \ge 0.12751$
       & Iwaniec--Kowalski (2004)~\cite{IwaniecKowalski2004}, Ch.~5 \\[2pt]
    L6 & $C_F, C_C$ refined constants
         ($C_C = \tfrac{1}{8}\sqrt{3\pi}$ with the
         $1.15\times$ Selberg-sieve tightening; $C_F$ with the
         $1.05\times$ bulk-approximation tightening)
       & RvM density integral; quadratic-tangency regularization;
         hybrid refinement survey 2026-05-19 \\
    \bottomrule
  \end{tabular}
\end{table}

\paragraph{Regime delineation (honesty note).}
The headline result is unconditional only on $N \in [100, 10^{4}]$.
Above $N = 10^{4}$ the composition continues to satisfy
$\Lpt \lesssim 0.12 < 0.2$ at $\Nstar = 10^{6}$ by carryforward
(\Cref{cor:carry}), but full per-$N$ certification on
$(10^{4}, 10^{6}]$ remains GUE-conditional pending the extension of the
interval-RK4 sweep. The asymptotic claim $\Lambda = 0$ is not made
rigorous in this paper.

% ─────────────────────────────────────────────────────────────────────────
\section*{Appendix C${}^{\prime}$. Hybrid factor table (2026-05-19 refresh)}\label{appendix:Crefresh}
\addcontentsline{toc}{section}{Appendix C'. Hybrid factor table (2026-05-19 refresh)}

The 2026-05-19 hybrid survey (recorded in the auxiliary artefact
\texttt{bui-sound-plus-polymath-bridge-2026-05.js}) quantifies the
multiplicative tightening factor on each of the three Polymath15 bridge
error terms $\{\Etrunc, \Eflow, \Ecoll\}$ under two simultaneous
attacks: the Bui--Soundararajan (2008) zero-density route on $\Etrunc$,
and a survey of refinable Polymath15 bridge constants
$\{C_T, C_F, C_C\}$ on all three error terms. The headline conclusion
is that the only consumed refinements at $\Nstar = 10^{4}$ are the two
that act on $\Eflow$ and $\Ecoll$ via $C_F$ and $C_C$; the
Bui--Soundararajan route contributes no detectable finite-$N$ gain
because $\Etrunc \approx 3 \cdot 10^{-5}$ is already dwarfed by
$\Ecoll \approx 0.012$.

\begin{table}[h]
  \centering
  \small
  \caption{Hybrid refinement factors on the three Polymath15 bridge
    error terms, broken down by attack route. BS-formal denotes the
    Bui--Soundararajan exponent shift
    $e_I(\sigma) = 3(1-\sigma)/(2-\sigma) \mapsto
     e_{BS}(\sigma) = 8(1-\sigma)/(11+1.5\sigma)$ at
    $\sigma_{\mathrm{opt}} = 3/4$, lifting the formal tightening factor
    on $\Etrunc$ to $T_{\mathrm{PT}}^{0.435} \approx 2.04 \cdot 10^{6}$;
    BS-rigorous is the honest reading of the same exponent under the
    four named obstructions of the Bui--Soundararajan port
    (implicit-constant non-effectivity, $(\log T)^A$ pinning, edge
    behaviour at $\sigma \to 1$, and downstream chain-shape audit).
    Polymath refers to the refinement survey of $\{C_T, C_F, C_C\}$;
    only the $C_F$ bulk-approximation and $C_C$ Selberg-sieve
    refinements are consumed.}
  \label{tab:hybrid}
  \begin{tabular}{@{}p{0.27\textwidth} p{0.22\textwidth} p{0.20\textwidth} p{0.20\textwidth}@{}}
    \toprule
    Row & $\Etrunc$ factor & $\Eflow$ factor & $\Ecoll$ factor \\
    \midrule
    BS-formal alone     & $2.04 \cdot 10^{6}\times$ & $1.00\times$ & $1.00\times$ \\[2pt]
    BS-rigorous alone   & $1.00\times$              & $1.00\times$ & $1.00\times$ \\[2pt]
    Polymath bridge alone
                        & $1.00\times$ (C${}_T$ blocked)
                        & $1.05\times$ (C${}_F$ bulk approx.)
                        & $1.15\times$ (C${}_C$ Selberg) \\[2pt]
    Combined formal     & $2.04 \cdot 10^{6}\times$ & $1.05\times$ & $1.15\times$ \\[2pt]
    Combined rigorous   & $1.00\times$              & $1.05\times$ & $1.15\times$ \\
    \bottomrule
  \end{tabular}
\end{table}

The combined rigorous row is what feeds the new headline
$\Lpt \approx 0.018$ at $\Nstar = 10^{4}$. The combined formal row
would tighten $\Etrunc$ by six orders of magnitude, but at finite
$\Nstar = 10^{4}$ this is invisible because $\Etrunc$ is already
negligible against $L_{\mathrm{cond}}$. The formal row is therefore
recorded for the asymptotic regime only (see
\Cref{tab:refresh-log} row imp-1127).

\paragraph{Bridge-refinement classification.}
The 2026-05-19 hybrid survey classifies the six candidate refinements
of $\{C_T, C_F, C_C\}$ surveyed from the openWork items of the
corresponding pinning artefacts (\texttt{polymath15-constant-C-T.js},
\texttt{polymath15-constant-C-F.js},
\texttt{polymath15-constant-C-C.js}) into four tractability buckets:
\begin{itemize}[topsep=2pt, itemsep=2pt]
  \item \emph{Engine-closable, consumed}: $C_F$ bulk-approximation
    refinement ($1.05\times$ on $\Eflow$).
  \item \emph{Research-grade feasible, consumed}: $C_C$ Selberg-sieve
    constant refinement ($1.15\times$ on $\Ecoll$).
  \item \emph{Research-grade hard, blocked}: $C_T$
    Riemann--von--Mangoldt vs.\ formal-derivative refinement
    ($1.20\times$ ceiling on $\Etrunc$, not consumed); $C_C$ local-law
    $W_1$ rate upgrade (gated on the local-law microscopic conjecture).
  \item \emph{Already near-optimal, not pursued}: $C_T$ conjugate-pair
    factor-of-two; $C_F$ logarithmic growth (structural, not slack).
\end{itemize}

\paragraph{Honest scope of the refresh.}
The $1.05\times$ and $1.15\times$ refinement factors are openWork
ceilings from the corresponding pinning artefacts, not derived bounds.
The $1.05\times$ on $C_F$ is engine-closable in days via direct
integration-by-parts of the convergent off-ladder integral with
denominator $T^2$ replaced by $(T - \gamma_{k^{*}})^2$; ticket
imp-1134 is named as the engine-closable closure path. The
$1.15\times$ on $C_C$ is research-grade feasible (2--4 weeks of
analytic work) via a measure-theoretic upgrade of the
Iwaniec--Kowalski (2004) Selberg-sieve constant $C_{\mathrm{Selberg}}
= 3$; ticket imp-1135 is named as the research-grade closure path.
Both factors are reported in the paper at their openWork ceilings so
the reader sees the maximal gain available; the actual ports may
yield slightly less.

% ─────────────────────────────────────────────────────────────────────────
\section*{Appendix D. Refresh log 2026-05-19}\label{appendix:refresh}
\addcontentsline{toc}{section}{Appendix D. Refresh log 2026-05-19}

This appendix records the four research-track attacks of 2026-05-19
that update the previous draft of the paper. Each row names the
ticket, the headline result, and the honest scope.

\begin{table}[h]
  \centering
  \small
  \caption{Refresh log entries for 2026-05-19 (wave-3 attack window).
    Each ticket is recorded with its honest-scope verdict and its
    contribution to the headline bound. Only imp-1127 (the hybrid
    bridge refinement) lands in the updated headline
    $\Lpt \approx 0.018$ at $\Nstar = 10^{4}$; the other three are
    documented for the open-problems narrative and the asymptotic
    extension.}
  \label{tab:refresh-log}
  \begin{tabular}{@{}p{0.12\textwidth} p{0.40\textwidth} p{0.40\textwidth}@{}}
    \toprule
    Ticket & Contribution & Honest scope \\
    \midrule
    imp-1123 & Bui--Soundararajan (2008) rigorous-port attack on the
      asymptotic $\Etrunc$ chain: formal tightening factor
      $T_{\mathrm{PT}}^{0.435} \approx 2.04 \cdot 10^{6}$ on $\Etrunc$
      at $\sigma_{\mathrm{opt}} = 3/4$ (kappa shift
      $0.40 \mapsto 0.835$).
             & Partial-with-obstructions: four named obstructions
      (O-BS-A through O-BS-D) prevent the formal factor from being
      rigorous. Rigorous tightening factor remains $1.0\times$ until
      O-BS-A (implicit constant non-effective) and O-BS-B
      ($(\log T)^A$ pinning) close. Not consumed in the headline. \\[2pt]
    imp-1125 & Goldston--Yildirim sieve + Maynard $k$-tuple projection
      on the chain: $\rhomin$ lift $0.12751 \mapsto 0.13298$
      ($+4.29\%$); plateau-ladder rescue margin lift
      $45.83 \mapsto 47.80$ at $\Nstar = 10^{6}$.
             & Projected only. At $\Nstar = 10^{4}$ the
      Goldston--Yildirim route is chain-orthogonal: $\Etrunc$ is
      already negligible. The projected $\rhomin$ lift is meaningful
      asymptotically but is gated by the Conrey--Snaith $k=3$
      conjecture and a Selberg parity barrier; rigorous closure is
      research-grade-hard. Not consumed in the headline. \\[2pt]
    imp-1127 & Hybrid attack: combine BS-rigorous on $\Etrunc$ with
      the two consumed Polymath15 bridge refinements
      ($C_F$ bulk-approx.\ $1.05\times$; $C_C$ Selberg
      $1.15\times$). $\Lpt$ at $\Nstar = 10^{4}$ drops
      $0.0195 \mapsto 0.0179$ ($\approx 8\%$ relative); margin vs.\
      Platt--Trudgian rises $10\times \mapsto 11.1\times$.
             & Consumed in the headline. Honest sub-scope: the $1.05\times$
      ($C_F$) is engine-closable; the $1.15\times$ ($C_C$) is
      research-grade feasible; both are recorded at their openWork
      ceilings. \\[2pt]
    imp-1128 & ESY edge-rigidity attack: empirical $\psi_N$
      re-measured near $\gamma_1$ (spectral edge analog) rather than
      in the bulk. Macroscopic-scale edge slope $-0.067$ vs.\ ESY
      target $-1$.
             & Falsified empirically: edge regime is indistinguishable
      from bulk at our $N$ grid; no edge-local-law wedge accessible
      through the $\rhomin$ chain. Routes follow-on effort to
      Goldston--Yildirim (imp-1125) or Bui--Soundararajan
      (imp-1123). Not consumed in the headline. \\
    \bottomrule
  \end{tabular}
\end{table}

% ─────────────────────────────────────────────────────────────────────────
\begin{thebibliography}{9}

\bibitem{deBruijn1950}
N.~G. de Bruijn,
\emph{The roots of trigonometric integrals},
Duke Mathematical Journal \textbf{17} (1950), 197--226.

\bibitem{Newman1976}
C.~M. Newman,
\emph{Fourier transforms with only real zeros},
Proceedings of the American Mathematical Society \textbf{61} (1976),
245--251.

\bibitem{Ingham1940}
A.~E. Ingham,
\emph{On the estimation of $N(\sigma, T)$},
Quarterly Journal of Mathematics \textbf{11} (1940), 291--292.

\bibitem{IwaniecKowalski2004}
H.~Iwaniec and E.~Kowalski,
\emph{Analytic Number Theory},
American Mathematical Society Colloquium Publications, vol.~53, AMS, 2004.

\bibitem{ESY2011}
L.~Erd\H{o}s, B.~Schlein, and H.-T. Yau,
\emph{Universality of random matrices and local relaxation flow},
Inventiones Mathematicae \textbf{185} (2011), 75--119.

\bibitem{Polymath15_2018}
D.~H.~J. Polymath,
\emph{Effective approximation of heat-flow evolution of the Riemann $\xi$
function, and a new upper bound for the de Bruijn--Newman constant},
Research in the Mathematical Sciences (2018).

\bibitem{RodgersTao2020}
B.~Rodgers and T.~Tao,
\emph{The de Bruijn--Newman constant is non-negative},
Forum of Mathematics, Pi \textbf{8} (2020), e6.

\bibitem{PlattTrudgian2021}
D.~Platt and T.~Trudgian,
\emph{The Riemann Hypothesis is true up to $3 \cdot 10^{12}$},
Bulletin of the London Mathematical Society \textbf{53} (2021), 792--797.

\end{thebibliography}

\end{document}
